% =========================================================
% Computational Linear Algebra — Lengths and Dot Products
% (Strang §1.2 / Johnston §1.2)
% =========================================================

\section{Lengths and Dot Products}

\begin{tcolorbox}[colback=gray!6, colframe=gray!40, arc=2pt,
  left=6pt, right=6pt, top=4pt, bottom=4pt,
  title={\small\textbf{Lengths and Dot Products --- Quick Reference}},
  fonttitle=\small\bfseries]
\begin{tabular}{@{}p{0.28\textwidth}p{0.65\textwidth}@{}}
\textbf{Dot product}     & $\mathbf{v} \cdot \mathbf{w} = v_1w_1 + v_2w_2 + \cdots + v_nw_n \in \mathbb{R}$. \\[4pt]
\textbf{Length (norm)}   & $\|\mathbf{v}\| = \sqrt{\mathbf{v} \cdot \mathbf{v}} = \sqrt{v_1^2 + \cdots + v_n^2}$. \\[4pt]
\textbf{Unit vector}     & $\hat{\mathbf{v}} = \mathbf{v}/\|\mathbf{v}\|$; has length 1. \\[4pt]
\textbf{Angle}           & $\cos\theta = \dfrac{\mathbf{v} \cdot \mathbf{w}}{\|\mathbf{v}\|\,\|\mathbf{w}\|}$. \\[6pt]
\textbf{Orthogonality}   & $\mathbf{v} \perp \mathbf{w}$ iff $\mathbf{v} \cdot \mathbf{w} = 0$. \\[4pt]
\textbf{Cauchy--Schwarz} & $|\mathbf{v} \cdot \mathbf{w}| \leq \|\mathbf{v}\|\,\|\mathbf{w}\|$. \\
\end{tabular}
\end{tcolorbox}

% ---------------------------------------------------------
\subsection{The Dot Product}

\begin{tcolorbox}[colback=propbox, colframe=propborder, arc=2pt,
  left=6pt, right=6pt, top=4pt, bottom=4pt,
  title={\small\textbf{Definition (Dot Product)}},
  fonttitle=\small\bfseries]
The \emph{dot product} (or \emph{inner product}) of $\mathbf{v}, \mathbf{w} \in \mathbb{R}^n$ is
\[
  \mathbf{v} \cdot \mathbf{w} = v_1 w_1 + v_2 w_2 + \cdots + v_n w_n \;\in\; \mathbb{R}.
\]
\end{tcolorbox}

\begin{tcolorbox}[colback=thmbox, colframe=thmborder, arc=2pt,
  left=6pt, right=6pt, top=4pt, bottom=4pt,
  title={\small\textbf{Theorem (Properties of the Dot Product)}},
  fonttitle=\small\bfseries]
For all $\mathbf{v}, \mathbf{w}, \mathbf{x} \in \mathbb{R}^n$ and $c \in \mathbb{R}$:
\begin{enumerate}
  \item $\mathbf{v} \cdot \mathbf{w} = \mathbf{w} \cdot \mathbf{v}$ \hfill (commutativity)
  \item $\mathbf{v} \cdot (\mathbf{w} + \mathbf{x}) = \mathbf{v} \cdot \mathbf{w} + \mathbf{v} \cdot \mathbf{x}$ \hfill (distributivity)
  \item $\mathbf{v} \cdot (c\mathbf{w}) = c(\mathbf{v} \cdot \mathbf{w})$ \hfill (scalar compatibility)
  \item $\mathbf{v} \cdot \mathbf{v} \geq 0$, with equality iff $\mathbf{v} = \mathbf{0}$ \hfill (positive definiteness)
\end{enumerate}
\end{tcolorbox}

% ---------------------------------------------------------
\subsection{Length and Angle}

\begin{tcolorbox}[colback=propbox, colframe=propborder, arc=2pt,
  left=6pt, right=6pt, top=4pt, bottom=4pt,
  title={\small\textbf{Definition (Length and Unit Vector)}},
  fonttitle=\small\bfseries]
The \emph{length} (or \emph{norm}) of $\mathbf{v} \in \mathbb{R}^n$ is
\[
  \|\mathbf{v}\| = \sqrt{\mathbf{v} \cdot \mathbf{v}} = \sqrt{v_1^2 + v_2^2 + \cdots + v_n^2}.
\]
A \emph{unit vector} has length $1$.
The \emph{normalisation} of a nonzero vector $\mathbf{v}$ is
$\hat{\mathbf{v}} = \mathbf{v}/\|\mathbf{v}\|$, which is the unit vector in the direction of $\mathbf{v}$.
\end{tcolorbox}

\begin{tcolorbox}[colback=thmbox, colframe=thmborder, arc=2pt,
  left=6pt, right=6pt, top=4pt, bottom=4pt,
  title={\small\textbf{Theorem (Cauchy--Schwarz Inequality)}},
  fonttitle=\small\bfseries]
For all $\mathbf{v}, \mathbf{w} \in \mathbb{R}^n$:
\[
  |\mathbf{v} \cdot \mathbf{w}| \leq \|\mathbf{v}\|\,\|\mathbf{w}\|,
\]
with equality if and only if $\mathbf{v}$ and $\mathbf{w}$ are parallel
(one is a scalar multiple of the other).
\end{tcolorbox}

\begin{tcolorbox}[colback=propbox, colframe=propborder, arc=2pt,
  left=6pt, right=6pt, top=4pt, bottom=4pt,
  title={\small\textbf{Definition (Angle Between Vectors)}},
  fonttitle=\small\bfseries]
The \emph{angle} $\theta \in [0, \pi]$ between nonzero vectors $\mathbf{v}, \mathbf{w} \in \mathbb{R}^n$
is defined by
\[
  \cos\theta = \frac{\mathbf{v} \cdot \mathbf{w}}{\|\mathbf{v}\|\,\|\mathbf{w}\|}.
\]
The vectors are \emph{orthogonal} (written $\mathbf{v} \perp \mathbf{w}$) if $\theta = \pi/2$,
equivalently if $\mathbf{v} \cdot \mathbf{w} = 0$.
\end{tcolorbox}

\begin{remark*}[Dot product and normals in NURBS]
The unit normal to a parametric surface $S(u,v)$ is computed as
$N = S_u \times S_v / \|S_u \times S_v\|$, which requires both the
cross product and normalisation via the norm.
The dot product also enters the convex hull argument: showing that a
B-spline curve lies on a specific side of a hyperplane reduces to
checking the sign of dot products of control points with the hyperplane's normal.
\end{remark*}
